\chapter{Connection Modeling with \texttt{port} and \texttt{connection}}

\begin{lernziele}
After this chapter, you can describe internal connections between system parts in SysML v2. You understand \texttt{port}, \texttt{connection}, and data flows, and you can carry these concepts over to ROS2-style communication in the robot project.
\end{lernziele}

\section{Introduction}
The structural modeling from the previous chapter shows which system parts exist. Connection modeling shows how these parts work together. It is the next step from structure to collaboration.

In SysML v2, you model connections with \texttt{port} and \texttt{connection}. Ports describe a system part's interfaces to the outside. Connections link ports of different parts to each other.

\begin{hinweisbox}
SysML v1 used the Internal Block Diagram (IBD) for this view. In SysML v2, ports and connections are direct model elements inside \texttt{part def}. The concept stays similar; the textual notation is more precise.
\end{hinweisbox}

\section{Purpose of Connection Modeling}
A connection view answers questions such as:

\begin{itemize}
    \item Which parts work together inside a system?
    \item Which data or signals flow between them?
    \item Which ports do the parts have?
    \item Which interfaces must be described?
    \item Where do integration risks arise?
\end{itemize}

\section{Ports}
A port describes an input or output of a system part. Ports make interfaces explicit and force you to name data or signal flows.

\begin{lstlisting}
port def ImagePort;
port def ScanPort;
port def ObjectListPort;

part def Camera {
  port imageOut : ImagePort;
}

part def Lidar {
  port scanOut : ScanPort;
}

part def ObjectDetector {
  port imageIn  : ~ImagePort;
  port objectsOut : ObjectListPort;
}

part def ObstacleDetector {
  port scanIn : ~ScanPort;
}
\end{lstlisting}

The tilde \texttt{\textasciitilde} marks a \glsdisp{conjport}{conjugated port}, which receives instead of sends.

\section{Connections}
\texttt{connection} links compatible ports to each other. A connection should always document what flows across it. The basic form is: \texttt{connection <name> connect <part1>.<port1> to <part2>.<port2>;}

\begin{lstlisting}
part def PerceptionSystem {
  part camera          : Camera;
  part lidar           : Lidar;
  part objectDetector  : ObjectDetector;
  part obstacleDetector : ObstacleDetector;

  connection c1 connect camera.imageOut
                      to objectDetector.imageIn;
  connection c2 connect lidar.scanOut
                      to obstacleDetector.scanIn;
}
\end{lstlisting}

Figure~\ref{fig:ports-connections} shows the resulting connection model with the ports involved.

\begin{figure}[H]
  \centering
  \adjustbox{max width=\linewidth,max totalheight=0.6\textheight}{% =====================================================================
% Figure: Connection model with ports "part def PerceptionSystem"
% Belongs to chapter 9.4/9.6 (listings "port def ...", "part def
% PerceptionSystem { ... connection c1 ... connection c2 ... }")
%
% Used in the book, e.g. in chapters_en/09_ibd.tex:
%
%   \begin{figure}[H]
%     \centering
%     % =====================================================================
% Figure: Connection model with ports "part def PerceptionSystem"
% Belongs to chapter 9.4/9.6 (listings "port def ...", "part def
% PerceptionSystem { ... connection c1 ... connection c2 ... }")
%
% Used in the book, e.g. in chapters_en/09_ibd.tex:
%
%   \begin{figure}[H]
%     \centering
%     % =====================================================================
% Figure: Connection model with ports "part def PerceptionSystem"
% Belongs to chapter 9.4/9.6 (listings "port def ...", "part def
% PerceptionSystem { ... connection c1 ... connection c2 ... }")
%
% Used in the book, e.g. in chapters_en/09_ibd.tex:
%
%   \begin{figure}[H]
%     \centering
%     \input{figures_en/fig_ports_connections}
%     \caption{Connection model of the \texttt{PerceptionSystem}. The small
%       squares on the box edges are ports. The tilde marks a conjugated
%       port, which receives instead of sends.}
%     \label{fig:ports-connections}
%   \end{figure}
%
% Required packages (already loaded in main_en.tex): tikz
% Required TikZ libraries: positioning, arrows.meta, shadows, calc
% =====================================================================

\adjustbox{max width=\textwidth,max totalheight=15.5cm}{\input{figures_en/fig_ports_connections.tikz}}

%     \caption{Connection model of the \texttt{PerceptionSystem}. The small
%       squares on the box edges are ports. The tilde marks a conjugated
%       port, which receives instead of sends.}
%     \label{fig:ports-connections}
%   \end{figure}
%
% Required packages (already loaded in main_en.tex): tikz
% Required TikZ libraries: positioning, arrows.meta, shadows, calc
% =====================================================================

\adjustbox{max width=\textwidth,max totalheight=15.5cm}{\begin{tikzpicture}[
    every node/.style={font=\small},
    part/.style={
        rectangle, rounded corners=1.5pt, draw=black!70, thick,
        minimum width=30mm, minimum height=12mm,
        align=center, fill=blue!8, drop shadow
    },
    port/.style={rectangle, draw=black, fill=black!75, minimum width=2.2mm,
        minimum height=2.2mm, inner sep=0pt},
    portlbl/.style={font=\small},
    flow/.style={-{Stealth[length=2.2mm]}, thick, orange!70!black}
]

% --- System parts ---
\node[part] (camera) at (0, 3)   {\texttt{camera :}\\\texttt{Camera}};
\node[part] (lidar)  at (0, 0)   {\texttt{lidar :}\\\texttt{Lidar}};
\node[part] (detect) at (7, 3)   {\texttt{objectDetector :}\\\texttt{ObjectDetector}};
\node[part] (obst)   at (7, 0)   {\texttt{obstacleDetector :}\\\texttt{ObstacleDetector}};

% --- Ports (small squares on the edge) ---
\node[port] (p-cam-out)   at ($(camera.east)$)   {};
\node[port] (p-lidar-out) at ($(lidar.east)$)    {};
\node[port] (p-det-in)    at ($(detect.west)$)   {};
\node[port] (p-det-out)   at ($(detect.east)$)   {};
\node[port] (p-obst-in)   at ($(obst.west)$)     {};

\node[portlbl, anchor=south west] at ([yshift=0.7mm]p-cam-out.north)   {\texttt{imageOut}};
\node[portlbl, anchor=south west] at ([yshift=0.7mm]p-lidar-out.north) {\texttt{scanOut}};
\node[portlbl, anchor=south east] at ([yshift=0.7mm]p-det-in.north)    {\texttt{imageIn} \textcolor{orange!70!black}{$\sim$}};
\node[portlbl, anchor=south west] at ([yshift=0.7mm]p-det-out.north)   {\texttt{objectsOut}};
\node[portlbl, anchor=south east] at ([yshift=0.7mm]p-obst-in.north)   {\texttt{scanIn} \textcolor{orange!70!black}{$\sim$}};

% --- Connections (connection c1 / c2) ---
\draw[flow] (p-cam-out)   -- node[above, midway, font=\small]{c1} (p-det-in);
\draw[flow] (p-lidar-out) -- node[above, midway, font=\small]{c2} (p-obst-in);

% --- Forwarding to navigation (objectsOut, not connected in the listing) ---
\draw[flow, dashed] (p-det-out) -- ++(2.2,0)
    node[right, font=\small, align=left] {to\\Navigation};

% --- Legend ---
\node[port, below=16mm of lidar, xshift=-4mm] (leg1) {};
\node[portlbl, right=1mm of leg1] {Port};
\node[portlbl, below=2mm of leg1, xshift=6mm] (legtilde) {\textcolor{orange!70!black}{$\sim$} = conjugated port (receives instead of sends)};
\draw[flow] ([yshift=-8mm]leg1.south) -- ++(6mm,0)
    node[right=1mm, font=\small] {Connection (\texttt{connection})};

\end{tikzpicture}
}

%     \caption{Connection model of the \texttt{PerceptionSystem}. The small
%       squares on the box edges are ports. The tilde marks a conjugated
%       port, which receives instead of sends.}
%     \label{fig:ports-connections}
%   \end{figure}
%
% Required packages (already loaded in main_en.tex): tikz
% Required TikZ libraries: positioning, arrows.meta, shadows, calc
% =====================================================================

\adjustbox{max width=\textwidth,max totalheight=15.5cm}{\begin{tikzpicture}[
    every node/.style={font=\small},
    part/.style={
        rectangle, rounded corners=1.5pt, draw=black!70, thick,
        minimum width=30mm, minimum height=12mm,
        align=center, fill=blue!8, drop shadow
    },
    port/.style={rectangle, draw=black, fill=black!75, minimum width=2.2mm,
        minimum height=2.2mm, inner sep=0pt},
    portlbl/.style={font=\small},
    flow/.style={-{Stealth[length=2.2mm]}, thick, orange!70!black}
]

% --- System parts ---
\node[part] (camera) at (0, 3)   {\texttt{camera :}\\\texttt{Camera}};
\node[part] (lidar)  at (0, 0)   {\texttt{lidar :}\\\texttt{Lidar}};
\node[part] (detect) at (7, 3)   {\texttt{objectDetector :}\\\texttt{ObjectDetector}};
\node[part] (obst)   at (7, 0)   {\texttt{obstacleDetector :}\\\texttt{ObstacleDetector}};

% --- Ports (small squares on the edge) ---
\node[port] (p-cam-out)   at ($(camera.east)$)   {};
\node[port] (p-lidar-out) at ($(lidar.east)$)    {};
\node[port] (p-det-in)    at ($(detect.west)$)   {};
\node[port] (p-det-out)   at ($(detect.east)$)   {};
\node[port] (p-obst-in)   at ($(obst.west)$)     {};

\node[portlbl, anchor=south west] at ([yshift=0.7mm]p-cam-out.north)   {\texttt{imageOut}};
\node[portlbl, anchor=south west] at ([yshift=0.7mm]p-lidar-out.north) {\texttt{scanOut}};
\node[portlbl, anchor=south east] at ([yshift=0.7mm]p-det-in.north)    {\texttt{imageIn} \textcolor{orange!70!black}{$\sim$}};
\node[portlbl, anchor=south west] at ([yshift=0.7mm]p-det-out.north)   {\texttt{objectsOut}};
\node[portlbl, anchor=south east] at ([yshift=0.7mm]p-obst-in.north)   {\texttt{scanIn} \textcolor{orange!70!black}{$\sim$}};

% --- Connections (connection c1 / c2) ---
\draw[flow] (p-cam-out)   -- node[above, midway, font=\small]{c1} (p-det-in);
\draw[flow] (p-lidar-out) -- node[above, midway, font=\small]{c2} (p-obst-in);

% --- Forwarding to navigation (objectsOut, not connected in the listing) ---
\draw[flow, dashed] (p-det-out) -- ++(2.2,0)
    node[right, font=\small, align=left] {to\\Navigation};

% --- Legend ---
\node[port, below=16mm of lidar, xshift=-4mm] (leg1) {};
\node[portlbl, right=1mm of leg1] {Port};
\node[portlbl, below=2mm of leg1, xshift=6mm] (legtilde) {\textcolor{orange!70!black}{$\sim$} = conjugated port (receives instead of sends)};
\draw[flow] ([yshift=-8mm]leg1.south) -- ++(6mm,0)
    node[right=1mm, font=\small] {Connection (\texttt{connection})};

\end{tikzpicture}
}
}
  \caption{Connection model of the \texttt{PerceptionSystem}. The small
    squares on the box edges are ports. The tilde marks a conjugated
    port, which receives instead of sends.}
  \label{fig:ports-connections}
\end{figure}

\section{Data Flows and ROS2}
Connection modeling fits ROS2 communication well. You can model ROS2 topics as data flows. A topic \texttt{/camera/image\_raw} connects \texttt{camera\_node} and \texttt{object\_detector\_node}. A topic \texttt{/cmd\_vel} connects navigation and motion control.

\begin{praxisbox}
A connection view for the \texttt{PerceptionSystem} shows: the camera delivers image data to the \texttt{ObjectDetector}, while the lidar delivers laser scan data to the \texttt{ObstacleDetector}. Both results go on to navigation. You can later derive ROS2 topics directly from this.
\end{praxisbox}

\section{A Connection Model for Navigation}
A connection view for navigation makes it visible that localization, map, path planning, and motion control work together:

\begin{lstlisting}
part def NavigationSystem {
  part localization : Localization;
  part mapProvider  : MapProvider;
  part pathPlanner  : PathPlanner;
  part pathFollower : PathFollower;

  connection c1 connect localization.poseOut
                      to pathPlanner.poseIn;
  connection c2 connect mapProvider.mapOut
                      to pathPlanner.mapIn;
  connection c3 connect pathPlanner.pathOut
                      to pathFollower.pathIn;
}
\end{lstlisting}

Here, \texttt{c1} delivers the current pose from localization to the path planner, \texttt{c2} delivers the map to that same path planner, and \texttt{c3} passes the planned path on to motion control.

\begin{figure}[H]
  \centering
  \adjustbox{max width=\linewidth,max totalheight=0.6\textheight}{% =====================================================================
% Figure: Connection model of the NavigationSystem
% Belongs to chapter 9, section "Connection Model for Navigation"
% (listing "part def NavigationSystem { ... connection c1/c2/c3 ... }")
%
% Used in the book, e.g. in chapters_en/09_ibd.tex:
%
%   \begin{figure}[H]
%     \centering
%     % =====================================================================
% Figure: Connection model of the NavigationSystem
% Belongs to chapter 9, section "Connection Model for Navigation"
% (listing "part def NavigationSystem { ... connection c1/c2/c3 ... }")
%
% Used in the book, e.g. in chapters_en/09_ibd.tex:
%
%   \begin{figure}[H]
%     \centering
%     % =====================================================================
% Figure: Connection model of the NavigationSystem
% Belongs to chapter 9, section "Connection Model for Navigation"
% (listing "part def NavigationSystem { ... connection c1/c2/c3 ... }")
%
% Used in the book, e.g. in chapters_en/09_ibd.tex:
%
%   \begin{figure}[H]
%     \centering
%     \input{figures_en/fig_kap09_verbindungsmodell_navigation}
%     \caption{Connection model of the \texttt{NavigationSystem}: localization
%       and map feed the path planner, whose path is passed on to motion
%       control.}
%     \label{fig:kap09-verbindungsmodell-navigation}
%   \end{figure}
%
% Required packages (already loaded in main_en.tex): tikz
% Required TikZ libraries: positioning, arrows.meta, shadows, calc
% =====================================================================

\adjustbox{max width=\textwidth,max totalheight=15.5cm}{\input{figures_en/fig_kap09_verbindungsmodell_navigation.tikz}}

%     \caption{Connection model of the \texttt{NavigationSystem}: localization
%       and map feed the path planner, whose path is passed on to motion
%       control.}
%     \label{fig:kap09-verbindungsmodell-navigation}
%   \end{figure}
%
% Required packages (already loaded in main_en.tex): tikz
% Required TikZ libraries: positioning, arrows.meta, shadows, calc
% =====================================================================

\adjustbox{max width=\textwidth,max totalheight=15.5cm}{\begin{tikzpicture}[
    every node/.style={font=\small},
    part/.style={
        rectangle, rounded corners=1.5pt, draw=green!40!black, thick,
        minimum width=32mm, minimum height=12mm,
        align=center, fill=green!15, drop shadow
    },
    port/.style={rectangle, draw=black, fill=black!75, minimum width=2.2mm,
        minimum height=2.2mm, inner sep=0pt},
    portlbl/.style={font=\small},
    flow/.style={-{Stealth[length=2.2mm]}, thick, orange!70!black}
]

% --- System parts ---
\node[part] (loc)  at (0, 3)   {\texttt{localization :}\\\texttt{Localization}};
\node[part] (map)  at (0, 0)   {\texttt{mapProvider :}\\\texttt{MapProvider}};
\node[part] (plan) at (6, 1.5) {\texttt{pathPlanner :}\\\texttt{PathPlanner}};
\node[part] (foll) at (11.5, 1.5) {\texttt{pathFollower :}\\\texttt{PathFollower}};

% --- Ports (small squares on the edge) ---
\node[port] (p-loc-out)  at ($(loc.east)$)  {};
\node[port] (p-map-out)  at ($(map.east)$)  {};
\node[port] (p-plan-posein) at ($(plan.west)!0.33!(plan.north west)$) {};
\node[port] (p-plan-mapin)  at ($(plan.west)!0.33!(plan.south west)$) {};
\node[port] (p-plan-out) at ($(plan.east)$) {};
\node[port] (p-foll-in)  at ($(foll.west)$) {};

\node[portlbl, anchor=south west] at ([yshift=0.7mm]p-loc-out.north)  {\texttt{poseOut}};
\node[portlbl, anchor=south west] at ([yshift=0.7mm]p-map-out.north)  {\texttt{mapOut}};
\node[portlbl, anchor=south east] at ([yshift=0.7mm]p-plan-posein.north) {\texttt{poseIn}};
\node[portlbl, anchor=north east] at ([yshift=-0.7mm]p-plan-mapin.south)  {\texttt{mapIn}};
\node[portlbl, anchor=south west] at ([yshift=0.7mm]p-plan-out.north) {\texttt{pathOut}};
\node[portlbl, anchor=north east] at ([yshift=-0.7mm]p-foll-in.south)  {\texttt{pathIn}};

% --- Connections (connection c1 / c2 / c3) ---
\draw[flow] (p-loc-out) -- node[above, midway, font=\small, sloped]{c1} (p-plan-posein);
\draw[flow] (p-map-out) -- node[below, midway, font=\small, sloped]{c2} (p-plan-mapin);
\draw[flow] (p-plan-out) -- node[above, midway, font=\small]{c3} (p-foll-in);

% --- Legend ---
\node[port, below=14mm of map, xshift=-4mm] (leg1) {};
\node[portlbl, right=1mm of leg1] {Port};
\draw[flow] ([yshift=-8mm]leg1.south) -- ++(6mm,0)
    node[right=1mm, font=\small] {Connection (\texttt{connection})};

\end{tikzpicture}
}

%     \caption{Connection model of the \texttt{NavigationSystem}: localization
%       and map feed the path planner, whose path is passed on to motion
%       control.}
%     \label{fig:kap09-verbindungsmodell-navigation}
%   \end{figure}
%
% Required packages (already loaded in main_en.tex): tikz
% Required TikZ libraries: positioning, arrows.meta, shadows, calc
% =====================================================================

\adjustbox{max width=\textwidth,max totalheight=15.5cm}{\begin{tikzpicture}[
    every node/.style={font=\small},
    part/.style={
        rectangle, rounded corners=1.5pt, draw=green!40!black, thick,
        minimum width=32mm, minimum height=12mm,
        align=center, fill=green!15, drop shadow
    },
    port/.style={rectangle, draw=black, fill=black!75, minimum width=2.2mm,
        minimum height=2.2mm, inner sep=0pt},
    portlbl/.style={font=\small},
    flow/.style={-{Stealth[length=2.2mm]}, thick, orange!70!black}
]

% --- System parts ---
\node[part] (loc)  at (0, 3)   {\texttt{localization :}\\\texttt{Localization}};
\node[part] (map)  at (0, 0)   {\texttt{mapProvider :}\\\texttt{MapProvider}};
\node[part] (plan) at (6, 1.5) {\texttt{pathPlanner :}\\\texttt{PathPlanner}};
\node[part] (foll) at (11.5, 1.5) {\texttt{pathFollower :}\\\texttt{PathFollower}};

% --- Ports (small squares on the edge) ---
\node[port] (p-loc-out)  at ($(loc.east)$)  {};
\node[port] (p-map-out)  at ($(map.east)$)  {};
\node[port] (p-plan-posein) at ($(plan.west)!0.33!(plan.north west)$) {};
\node[port] (p-plan-mapin)  at ($(plan.west)!0.33!(plan.south west)$) {};
\node[port] (p-plan-out) at ($(plan.east)$) {};
\node[port] (p-foll-in)  at ($(foll.west)$) {};

\node[portlbl, anchor=south west] at ([yshift=0.7mm]p-loc-out.north)  {\texttt{poseOut}};
\node[portlbl, anchor=south west] at ([yshift=0.7mm]p-map-out.north)  {\texttt{mapOut}};
\node[portlbl, anchor=south east] at ([yshift=0.7mm]p-plan-posein.north) {\texttt{poseIn}};
\node[portlbl, anchor=north east] at ([yshift=-0.7mm]p-plan-mapin.south)  {\texttt{mapIn}};
\node[portlbl, anchor=south west] at ([yshift=0.7mm]p-plan-out.north) {\texttt{pathOut}};
\node[portlbl, anchor=north east] at ([yshift=-0.7mm]p-foll-in.south)  {\texttt{pathIn}};

% --- Connections (connection c1 / c2 / c3) ---
\draw[flow] (p-loc-out) -- node[above, midway, font=\small, sloped]{c1} (p-plan-posein);
\draw[flow] (p-map-out) -- node[below, midway, font=\small, sloped]{c2} (p-plan-mapin);
\draw[flow] (p-plan-out) -- node[above, midway, font=\small]{c3} (p-foll-in);

% --- Legend ---
\node[port, below=14mm of map, xshift=-4mm] (leg1) {};
\node[portlbl, right=1mm of leg1] {Port};
\draw[flow] ([yshift=-8mm]leg1.south) -- ++(6mm,0)
    node[right=1mm, font=\small] {Connection (\texttt{connection})};

\end{tikzpicture}
}
}
  \caption{Connection model of the \texttt{NavigationSystem}: localization and map feed the path planner, whose path is passed on to motion control.}
  \label{fig:kap09-verbindungsmodell-navigation}
\end{figure}

Navigation emerges from several parts working together (see figure~\ref{fig:kap09-verbindungsmodell-navigation}). At the same time, this makes visible the interfaces that you can later implement in ROS2.

\section{Documenting Assumptions}
A good connection model should also document assumptions. Is the map loaded beforehand, or updated while driving? Does localization work relative to \texttt{map} or \texttt{odom} (two common ROS2 coordinate frames, explained in more detail in chapter 10)? Such questions belong in the model element as documentation (\texttt{doc /* ... */}).

\begin{lstlisting}
part def NavigationSystem {
  doc /* Precondition: an initial map is available.
         Localization works relative to the map frame. */
  ...
}
\end{lstlisting}

\section{Common Mistakes}
\subsection{Connecting Parts Without Ports}
Connections should always run through named ports. Direct connections without ports leave open what exactly gets transmitted.

\subsection{Ports Without Data Types}
A port only becomes useful once you know what data flows across it. Define \texttt{port def} with descriptive names.

\subsection{Too Much Detail in One Model Section}
A connection view should answer one clear question. For large systems, several focused sections work better than one overloaded model.

\section{Summary}
Connection modeling with \texttt{port} and \texttt{connection} shows how system parts work together internally. It is especially useful for interfaces and data flows. In the robot project, it bridges the SysML v2 structure and ROS2 communication.

\section{Review Questions}
\begin{enumerate}
    \item What does a connection view show that a structural view does not?
    \item What does a \texttt{port} describe?
    \item What does the tilde do on a port?
    \item What is \texttt{connection} used for?
    \item Why is connection modeling helpful for ROS2?
    \item Which topic could sit between the camera and the ObjectDetector?
    \item Why should a connection model document its assumptions?
\end{enumerate}

\section{Exercises}
\begin{uebungbox}
\textbf{Exercise 1: Connection model for PerceptionSystem}\\
Create a \texttt{part def PerceptionSystem} with a camera, lidar, ObjectDetector, and ObstacleDetector. Define ports and connect them with \texttt{connection}.
\end{uebungbox}

\begin{uebungbox}
\textbf{Exercise 2: Name ports}\\
Define at least two ports with descriptive names for each of \texttt{BatterySystem}, \texttt{NavigationSystem}, and \texttt{MotionControl}.
\end{uebungbox}

\begin{uebungbox}
\textbf{Exercise 3: Prepare ROS2 mapping}\\
Pick three connections from your model and assign each one a possible ROS2 topic.
\end{uebungbox}

\section{Further Reading}
Ports and connections in SysML v2 are documented in the SysML v2 pilot project \cite{sysmlv2pilot}. The Syside Editor offers concrete examples of connection modeling \cite{syside}. For implementation in ROS2, the ROS2 documentation and relevant introductions are helpful \cite{ros2docs}.
